\chapter{总结与展望}\label{chap:conclusion}



\section{工作总结}
本文围绕后量子时代端到端加密通信协议的设计与实现，聚焦于解决现有方案通信开销大、国产化研究不足等关键问题，取得了一系列研究成果：

\begin{enumerate}
    \item \textbf{提出了一种基于 Tag-KEM 的高效后量子混合协议}：通过引入标签密钥封装机制（Tag-KEM），将后量子共享密钥与通信上下文（如临时公钥、身份信息）进行强绑定，成功解决了后量子密钥协商过程中密文复用与跨会话密钥隔离的难题。该设计允许同一对后量子公私钥安全地服务于多个不同会话，显著降低了带宽消耗和计算开销。

    \item \textbf{构建了完整的安全模型并进行了形式化证明}：在 Cohn-Gordon 等人工作的基础上，建立了适用于多阶段、混合模式密钥交换协议的安全模型。通过博弈序列（Game-hopping）方法，严格证明了所提协议在新鲜度条件下满足会话密钥不可区分性（ms-Ind），为协议的前向安全与后向安全提供了坚实的理论保障。

    \item \textbf{完成了协议的工程化实现与国产化适配}：在 \texttt{libsignal} 协议框架下，不仅实现了基于国际标准 ML-KEM（Kyber）的 TKEM 方案，还成功集成了我国自主设计的后量子算法 Aigis-Enc。同时，全面替换了协议中的核心密码组件，采用 SM2 进行密钥协商与签名、SM3 作为哈希与 KDF、SM4 进行对称加密，构建了一套完全自主可控的端到端加密通信体系。

    \item \textbf{设计了基于 TEE 的安全增强方案}：针对后量子私钥需长期存储带来的安全风险，设计并实现了基于 Intel SGX 的可信执行环境（TEE）安全架构。该方案将私钥的生成、存储和解封装全过程置于硬件保护的 Enclave 内，有效抵御了操作系统层面的内存窃取和离线攻击，显著提升了系统的整体安全性。

    \item \textbf{通过实验验证了方案的有效性}：性能评估结果表明，本文提出的方案在引入后量子安全性和国密算法后，其通信延迟和计算开销均在可接受范围内，证明了该方案在实际应用中的可行性与优越性。
\end{enumerate}

\section{研究展望}
尽管本研究取得了一定成果，但仍存在一些值得深入探索的方向：

\begin{enumerate}
    \item \textbf{群聊场景的扩展}：本文工作主要针对点对点通信场景。未来可将基于 Tag-KEM 的后量子安全机制扩展至群组通信协议中，研究如何在保证成员动态加入/退出、前向/后向安全的同时，高效管理群组的后量子密钥材料。

    \item \textbf{更多国产后量子算法的集成}：目前仅集成了 Aigis-Enc 作为国产后量子 KEM 的代表。未来可进一步适配如 CTRU/C-NTRU 等其他优秀的国产后量子密码方案，并对它们在协议中的性能和安全性进行横向对比评估。

    \item \textbf{跨平台 TEE 支持}：当前的 TEE 安全方案基于 Intel SGX，主要适用于 x86 架构。未来可探索在 ARM TrustZone、RISC-V Keystone 等不同硬件平台上实现类似的安全架构，以提升方案在移动设备和物联网等场景下的普适性。

    \item \textbf{标准化与互操作性研究}：推动所设计的混合模式协议及其国密算法套件的标准化进程，并与其他后量子通信协议（如 Open Quantum Safe 项目）进行互操作性测试，是实现大规模落地应用的关键一步。
\end{enumerate}